% ===================================================================
%  ТАБЛИЦЯ № 4 — Зрілість і Старість.
%  Traduction 2026 d'apres texto/io, controlee sur texto/fr, texto/en
%  et texto/es. Consigne : voir texto/uk/10-tablycia-01.tex.
%
%  LE TABLEAU DE LA PARENTE. Vingt-deux degres nommes en deux scenes,
%  et c'est le lexique ou les deux langues slaves se separent le plus
%  nettement : le garcon d'honneur est un « дружба », la demoiselle
%  d'honneur une « дружка », et la sœur du mari une « зовиця » — trois
%  mots que le russe ne connait pas sous cette forme. Les mois d'hiver
%  non plus ne se ressemblent en rien : « декабрь » est « грудень ».
% ===================================================================

%%K t04-apar-1 apar
\VUsaut{4.0mm}
\VUtitre{15.0pt}{60}{ТАБЛИЦЯ № 4}
\VUsaut{1.6mm}
\VUtitre{11.4pt}{0}{Зрілість і Старість.}
\VUsaut{3.2mm}

%%K t04-tit-1 sub
\VUcentre{11.4pt}{0}{\textit{Перша сцена.}}
\VUsaut{1.6mm}
\VUtitre{11.4pt}{0}{Весільний обід.}
\VUsaut{2.6mm}

%%K t04-tit-2 sub
\VUpk{10.2pt}{40}{Осінь.}
\VUsaut{3.0mm}

%%K t04-c1-01-1 p
1. --- Молоді люди виросли. Один із них, Іван, жениться. Ми на
\VUgras{весільному обіді}, який зібрав за одним столом родини нареченого і
нареченої.

%%K t04-c1-02-1 p
2. --- Стоїть \VUgras{осінь}, місяць \VUgras{вересень}. Холодний вітер
\VUgras{жовтня} і \VUgras{листопада} ще не зірвав із полів зеленого листя.
Вечірнє повітря м'яке й запашне; тому обід подають під \VUgras{верандою}
\textsuperscript{(8)}, що виходить у сад.

%%K t04-c1-03-1 p
3. --- Удалині, на \VUgras{пагорбі} \textsuperscript{(3)}, стоїть дім Іванових
батьків, зграбна \VUgras{садиба} \textsuperscript{(1)}, схована в зелені;
прекрасні виноградники, що дають добірне вино, вкривають схил пагорба.
Виноградар ходить коло них і плекає їх.

%%K t04-c1-04-1 p
4. --- Над лозами пролітає зграйка \VUgras{куріпок} \textsuperscript{(2)},
сполохана наближенням \VUgras{лісника} \textsuperscript{(5)}. А той, з
\VUgras{рушницею} під пахвою і \VUgras{ловецькою торбою} через плече, обшукує
\VUgras{кущі} \textsuperscript{(4)} зі своїм \VUgras{собакою}
\textsuperscript{(6)}. Він стереже маєток і не дає браконьєрам винищувати
дичину.

%%K t04-c1-05-1 p
5. --- Знадвору по \VUgras{веранді} в'ється \VUgras{шпалерна лоза}, з якої
звисають важкі \VUgras{грона} \textsuperscript{(7)}.

%%K t04-c1-06-1 p
6. --- Прекрасні осінні квіти, що прикрашають \VUgras{стіл}
\textsuperscript{(16)}, --- це \VUgras{хризантеми} \textsuperscript{(9)};
\VUgras{батько} \textsuperscript{(10)} нареченої встромив одну в петельку фрака.

%%K t04-c1-07-1 p
7. --- Обід доходить кінця. \VUgras{Слуга} \textsuperscript{(11)} починає вносити
страви десерту і скоро прибере вже непотрібні частини приборів ---
\VUgras{ложки} \textsuperscript{(14)}, \VUgras{виделки} \textsuperscript{(12)},
\VUgras{ножі} \textsuperscript{(13)}, \VUgras{полумиски}
\textsuperscript{(15)}.

%%K t04-tit-3 sub
\VUpk{10.2pt}{40}{Родина.}
\VUsaut{3.0mm}

%%K t04-c1-08-1 p
8. --- У всіх вигляд веселий, а усмішка, з якою \VUgras{мати}
\textsuperscript{(17)} нареченого дивиться на своїх дітей, показує, яка вона
щаслива.

%%K t04-c1-09-1 p
9. --- Цей шлюб єднає дві багаті родини округи. Іван, \VUgras{чоловік}
\textsuperscript{(18)}, --- \VUgras{син} великого землевласника, і він стає
\VUgras{зятем} заможного фабриканта.

%%K t04-c1-10-1 p
10. --- \VUgras{Подружжя} молоде, і все ж вони довго були \VUgras{заручені}.
\VUgras{Молода дружина} \textsuperscript{(19)}, Марта, мила у своїй \VUgras{білій
сукні}, прибраній флердоранжем.

%%K t04-c1-11-1 p
11. --- Її \VUgras{свекор} \textsuperscript{(20)} радий такій добрій
\VUgras{невістці}, а \VUgras{теща} \textsuperscript{(21)} Івана усміхається,
бачачи свою \VUgras{дочку} такою гарною.

%%K t04-c1-12-1 p
12. --- На кінці столу сидять решта \VUgras{гостей} \textsuperscript{(22)}:
\VUgras{рідня, друзі, двоюрідні}. \VUgras{Метрдотель} \textsuperscript{(23)} із
серветкою на руці прудко їм прислуговує.

%%K t04-tit-4 sub
\VUpk{10.2pt}{40}{Друзі.}
\VUsaut{3.0mm}

%%K t04-c1-13-1 p
13. --- По праву руч столу --- \VUgras{панночка} \textsuperscript{(24)},
\VUgras{подруга} нареченої, тоді \VUgras{брат} нареченої, її \VUgras{дружба}
\textsuperscript{(25)}. Він розмовляє зі своєю \VUgras{дружкою}
\textsuperscript{(26)}, сестрою нареченого, яка через цей шлюб стає
\VUgras{зовицею} Марти. Це чарівна молода жінка. У неї гарні
\VUgras{коштовності}, \VUgras{перлове намисто} і золотий \VUgras{браслет}.

%%K t04-c1-14-1 p
14. --- Поруч із ними пан, який устав і підняв келих, --- \VUgras{друг}
\textsuperscript{(27)} родини. Він один зі \VUgras{свідків нареченого}. Він каже
\VUgras{тост}, п'є за \VUgras{здоров'я} молодих і зичить їм довгого щастя. Він у
парадному вбранні: фрак і \VUgras{ляковані туфлі} \textsuperscript{(28)}. Він
веде під руку \VUgras{бабусю} \textsuperscript{(29)}, яка \VUgras{вдова}.

%%K t04-c1-15-1 p
15. --- Молодий чоловік у \VUgras{фраку} \textsuperscript{(31)}, з тонкими
чорними вусами, --- \VUgras{шурин} \textsuperscript{(30)} Івана. Він визначний
інженер. Він сидить поруч із дуже зграбною \VUgras{панією}
\textsuperscript{(33)}, \VUgras{старшою сестрою} Марти і матір'ю тієї милої
дівчинки, яка йде під руку з двоюрідним братом піднести квітку і
\VUgras{побажати} \VUgras{доброго вечора}. У цієї пані прекрасні
\VUgras{сережки} і зграбна \VUgras{зачіска}.

%%K t04-c1-16-1 p
16. --- Малий двоюрідний брат, такий кумедний у своїй \VUgras{блузі}
\textsuperscript{(36)} і широких \VUgras{штанцях} \textsuperscript{(37)}, вельми
гідний жалю. Він \VUgras{сирота} \textsuperscript{(35)}. Він утратив батька й
матір зовсім малим. Його дядько --- його \VUgras{опікун}
\textsuperscript{(38)} і зостався неодруженим, щоб виростити бідну дитину. Він
добра людина. Він добряче лисий і дуже короткозорий; він носить
\VUgras{пенсне}.

%%K t04-tit-5 sub
\VUsaut{2.4mm}
\VUpk{10.2pt}{40}{Десерт.}
\VUsaut{3.0mm}

%%K t04-c1-17-1 p
17. --- Позаду гостей, на столі, вкритому \VUgras{скатертиною}, розкладено
\VUgras{десерт} \textsuperscript{(39)}. У великій вазі --- садовина:
\VUgras{персики} \textsuperscript{(40)}, \VUgras{груші} \textsuperscript{(41)},
\VUgras{яблука} \textsuperscript{(42)}, \VUgras{абрикоси} \textsuperscript{(43)}
і \VUgras{виноград} \textsuperscript{(44)}. Поруч --- \VUgras{ананаси}
\textsuperscript{(45)}, прекрасний \VUgras{торт} \textsuperscript{(46)},
\VUgras{пляшки лікеру} \textsuperscript{(48)} і \VUgras{шампанського}
\textsuperscript{(49)}, а також стоси чистих \VUgras{тарілок}
\textsuperscript{(47)}.

%%K t04-c1-18-1 p
18. --- \VUgras{Чашник} \textsuperscript{(50)} відкорковує \VUgras{пляшку}
\textsuperscript{(52)} старого вина \VUgras{коркотягом} \textsuperscript{(51)}.
\VUgras{Корок} \textsuperscript{(53)} іде, здається, дуже туго. У цього чашника
на руці \VUgras{серветка} \textsuperscript{(54)}.

%%K t04-c1-19-1 p
19. --- Після обіду чоловіки перейдуть до курильні, а пані підуть одягатися до
балу.

%%K t04-tit-6 sub
\VUsaut{5.0mm}
\VUcentre{11.4pt}{0}{\textit{Друга сцена.}}
\VUsaut{1.6mm}
\VUpk{11.4pt}{40}{Дідусеві іменини.}
\VUsaut{2.6mm}

%%K t04-tit-7 sub
\VUpk{10.2pt}{40}{Відвідини.}
\VUsaut{3.0mm}

%%K t04-c2-01-1 p
1. --- Минуло десять років. Іванові батьки постаріли і ніжно люблять своїх
трьох онучат: двох \VUgras{онуків} \textsuperscript{(2)} і \VUgras{онуку}
\textsuperscript{(3)}. Оскільки сьогодні \VUgras{дідусеві іменини}, діти прийшли
його привітати.

%%K t04-c2-02-1 p
2. --- Малий Петро ще зовсім малий; йому тільки три роки. Він бачить, як до
нього йде \VUgras{кішка} \textsuperscript{(1)}. Йому хочеться погладити її гарну
шовковисту \VUgras{шерсть} і довгий \VUgras{хвіст}; йому й на думку не спадає,
що під цими зграбними \VUgras{лапками} є гострі злі пазурці, які можуть його
подряпати.

%%K t04-c2-03-1 p
3. --- Його сестра Габріель, яка тримає його за руку, куди поважніша, а Марсель
у великому \VUgras{пальті} \textsuperscript{(4)} і шкіряному \VUgras{поясі}
\textsuperscript{(5)} виглядає зовсім дорослим, попри короткі штани, з яких
видно його \VUgras{панчохи} \textsuperscript{(6)}.

%%K t04-c2-04-1 p
4. --- Їхній батько надів товсте зимове \VUgras{пальто} \textsuperscript{(7)} з
\VUgras{хутряним коміром} \textsuperscript{(8)}, а в \VUgras{кишені}
\textsuperscript{(9)} в нього шарф. На його дружині повстяний капелюх із
\VUgras{пір'ям} \textsuperscript{(10)}, \VUgras{жакет} \textsuperscript{(11)},
хутряне \VUgras{боа} \textsuperscript{(12)} і \VUgras{муфта}
\textsuperscript{(13)}.

%%K t04-c2-05-1 p
5. --- Дуже холодно, бо прийшла зима. \VUgras{Сніг} \textsuperscript{(19)} ішов
цілий день, і дахи домів зовсім білі. З \VUgras{ніччю} мороз дужчає ще більше. У
небі світять \VUgras{місяць} \textsuperscript{(17)} і \VUgras{зорі}
\textsuperscript{(18)}, пронизуючи \VUgras{темряву}.

%%K t04-tit-8 sub
\VUsaut{4.2mm}
\VUpk{10.2pt}{40}{Хвороба.}
\VUsaut{3.0mm}

%%K t04-c2-06-1 p
6. --- Бідний \VUgras{сліпий} \textsuperscript{(20)}, який переходить площу
перед \VUgras{аптекою} \textsuperscript{(16)}, ведений своїм \VUgras{пуделем}
\textsuperscript{(21)}, дуже нещасний. Юстина, \VUgras{служниця}
\textsuperscript{(14)}, несе з кухні \VUgras{чашку} \textsuperscript{(15)} дуже
гарячого \VUgras{трав'яного настою}, щедро підсолодженого.

%%K t04-c2-07-1 p
7. --- \VUgras{Бабуся} \textsuperscript{(23)}, якій хочеться взяти зі столу свої
\VUgras{окуляри} \textsuperscript{(26)}, щоб прочитати \VUgras{рецепт}
\textsuperscript{(27)}, щойно написаний \VUgras{лікарем} \textsuperscript{(25)},
підводиться з крісла, спираючись на свою \VUgras{палицю}
\textsuperscript{(24)}.

%%K t04-c2-08-1 p
8. --- Вона часто потерпає від дрібних \VUgras{недуг}, \VUgras{запаморочень},
\VUgras{застуд} і \VUgras{зубного болю}; ноги в неї слабкі, а очі вже
недобачають. І все ж вона рада подражнити свого старого \VUgras{лікаря}, який
завжди хоче приписати їй \VUgras{мікстуру} \textsuperscript{(28)}. Сьогодні вона
дуже щаслива, що молода рідня навколо неї.

%%K t04-c2-09-1 p
9. --- У добре зачиненій кімнаті затишно. У мармуровому \VUgras{каміні}
\textsuperscript{(29)} палає \VUgras{славний вогонь} \textsuperscript{(30)}, і
добре на душі при м'якому \VUgras{світлі} щойно засвіченої \VUgras{лампи}
\textsuperscript{(31)}.

%%K t04-tit-9 sub
\VUsaut{4.2mm}
\VUpk{10.2pt}{40}{Дідусь.}
\VUsaut{3.0mm}

%%K t04-c2-10-1 p
10. --- Навіть \VUgras{дідусь} \textsuperscript{(33)} забув, що був хворий, що
\VUgras{ревматизм} тримає його прикутим до \VUgras{крісла}
\textsuperscript{(34)} і що ще вчора в нього були \VUgras{жар і млості}. Він уже
не пам'ятає, що минулої зими в нього було \VUgras{запалення легенів} і що
жорстокий, болісний \VUgras{кашель} завдав йому чимало муки.

%%K t04-c2-10-2 p
А втім, оклигав він насилу. Тепер йому краще. Але нога, яку він тримає на
\VUgras{подушці} \textsuperscript{(38)}, покладеній на низький \VUgras{ослінчик}
\textsuperscript{(39)}, усе ще болить у нього.

%%K t04-c2-11-1 p
11. --- Коли він дбайливо закутаний у свій теплий \VUgras{халат}
\textsuperscript{(35)} з великими перламутровими \VUgras{ґудзиками}
\textsuperscript{(36)} і коли поруч із ним, щоб читати йому газету, сидить мала
\VUgras{сирітка} \textsuperscript{(40)} в білому \VUgras{чепчику}, в синій
\VUgras{сукні} \textsuperscript{(42)} і \VUgras{пелеринці}
\textsuperscript{(41)}, він не нарікає.

%%K t04-c2-12-1 p
12. --- Щороку, коли \VUgras{календар} \textsuperscript{(43)} знову приводить
\VUgras{число} першого \VUgras{грудня}, він з нетерпінням жде своїх
\VUgras{онуків}, бо знає, що вони не забудуть цього важливого \VUgras{числа}.

%%K t04-c2-13-1 p
13. --- Він зворушено згадує далекі дні, коли сам ходив вітати славного
імператорського \VUgras{маршала}, чий \VUgras{портрет} \textsuperscript{(44)}
висить на стіні і який доводиться дітям \VUgras{прадідом}.
\VUsaut{6.0mm}
\VUfilet{22mm}
